%% kapa: exam=1 round=30 year=2019 period=1 form=A subject=economics subject_ko=경제학원론 seq=5 number=45
%% source: Q-Net 감정평가사 기출문제 (cst00309) artlSeq=5205205 file=extracted/1008/5205205/2019년 제30회 감정평가사 1차시험 문제지/2019년 제30회 감정평가사 1차시험 문제지/2019년도 제30회 감정평가사 1차 1교시 A형.hwp
%% conversion: kordoc 4.24.51 -> tools/gen1.py (auto); review=proofread-sonnet; flags=none
\begin{question}{45}
기업의 생산기술이 진보하는 경우에 관한 설명으로 옳은 것을 모두 고른 것은?
\begin{framedtext}
ㄱ. 자본절약적 기술진보가 일어나면 평균비용곡선이 하방 이동한다. \par
ㄴ. 자본절약적 기술진보가 일어나면 등량곡선이 원점에서 멀어진다. \par
ㄷ. 노동절약적 기술진보가 일어나면 한계비용곡선이 하방 이동한다. \par
ㄹ. 중립적 기술진보가 일어나면 노동의 한계생산 대비 자본의 한계생산은 작아진다.
\end{framedtext}
\begin{choices}
  \item ㄱ, ㄴ
  \item ㄱ, ㄷ
  \item ㄴ, ㄷ
  \item ㄴ, ㄹ
  \item ㄷ, ㄹ
\end{choices}
\end{question}
\answer{2}
